\documentclass[11pt]{article}
\usepackage[margin=1in]{geometry}
\usepackage{amsmath,amssymb,amsthm}
\usepackage[colorlinks=true,urlcolor=blue,linkcolor=blue]{hyperref}
\newtheorem{theorem}{Theorem}
\newtheorem{lemma}[theorem]{Lemma}
\newtheorem{corollary}[theorem]{Corollary}
\newtheorem{definition}[theorem]{Definition}
\newcommand{\OPT}{\operatorname{OPT}}
\newcommand{\MS}{\mathrm{MS}}
\newcommand{\dep}{\operatorname{depth}}
\title{The final switch of multi-splay: a two-copy embedding}
\author{Research note; hand argument}
\date{2 October 2026}
\begin{document}
\maketitle

\noindent\textbf{Status.}
This note isolates an endpoint issue in comparing standard multi-splay
with the local Cha-Cha lower bound. Its main theorem is a hand
argument conditional on that manuscript's dynamic lower-bound theorem.
The copy embedding, comparator lifting, and initialization calculation
were independently hand-audited in this research session. The extension
to arbitrary depth hierarchies is a hand extension of the manuscript;
the existing Lean endpoint is stated for its midpoint hierarchy.
It does not prove the desired multi-splay upper bound. In particular,
the secondary counter below is an analytical ledger, not an assertion
about the physical order of the multi-splay operations.

\section{The counter to be bounded}

Let $P$ be the perfect node-oriented reference BST on
$n=2^h-1$ keys, whose node depths are $0,\ldots,h-1$.
The initial preferred direction is right at every node. A query $x$
performs the standard multi-splay switches at the ancestors whose
preferences point away from $x$, followed by the final switch at $x$.
The physical ancestor switches run from bottom to top. The analytical
emission list $E_i$ instead lists their depths in increasing order and
appends $\dep(x)$ for the final switch. Thus $E_i$ is strictly increasing.
Write $S_{\MS}=\sum_i|E_i|$.

For a leaf of $P$, the final operation still performs its splays although
there are no children to mark. Give such a leaf a fictitious direction
bit, toggled by this final operation. This bit has no effect on the
physical algorithm; it only makes the encoding uniform.

Let $B$ be any binary interval hierarchy on the depth labels
$0,\ldots,h$. For an internal interval $I=[a,b]$ split at $c$, and a
primary node $u$ of depth $a$, retain an analytical bit $b_{u,I}$,
initially right. An emission $d\in E_i\cap I$ writes left when $d\le c$
and right otherwise, to the bit owned by the depth-$a$ ancestor of
$x_i$. Process $E_i$ in increasing order. Let $F_{\MS}$ count changes of
these bits. There is no extra completion emission $h$ in $E_i$.

\begin{theorem}\label{thm:endpoint}
Assume the dynamic lower-bound theorem in the local Cha-Cha
manuscript. Then, for the perfect reference above,
\[
 |X|+S_{\MS}(X)+F_{\MS}(X)=O(\OPT(X)+n).
\]
The comparator may choose its initial BST freely. The constant is
independent of $B$ and $h$.
\end{theorem}

The theorem avoids the invalid shortcut of charging every final primary
switch only a constant number of secondary flips. A single extra
emission can change several secondary bits. The proof instead embeds
the entire endpoint-aware history into a legitimate leaf-query history.

\section{The lower bound permits arbitrary interval hierarchies}

The Cha-Cha manuscript defines a midpoint hierarchy, but its
dynamic lower-bound proof uses only the following properties of that
hierarchy:
\begin{enumerate}
\item For two distinct depth labels, exactly one split separates them.
This is their least common ancestor in the hierarchy.
\item The two children of an interval $[a,b]$ are $[a,c]$ and
$[c+1,b]$. The macro word of the parent is the first child's word;
the micro words are the second child's words, over its different owners.
\item Every nonroot fragment has exactly one parent, so the turn
potentials telescope to the single-depth leaves.
\end{enumerate}
The subsequence-monotonicity proof and the injection of entry and delayed
witnesses use the first property. The normalized MTR simulation and
return property do not use secondary split positions. The turn identity
and its finite occurrence accounting use interval partitions, and the
final telescoping uses the last two properties. No comparison between
the lengths of the two children occurs in these arguments.
Consequently the same proof gives, for any binary interval hierarchy,
\begin{equation}\label{eq:lower}
 |Z|+P(Z)+D_B(Z)=O(\OPT(Z)).
\end{equation}
Balancing is relevant to the upper-bound implementation, but not to this
lower-bound extension.

\section{A perfect leaf reference with two copies of each key}

Build a full binary primary reference $T$ of height $2h$. An original
node $q$ at depth $d$ corresponds to a pivot $v(q)$ at depth $2d$.
Its two children are connector nodes. Their ordered children are
\[
 \bigl(\text{recursive left subtree of }q,\ q^-\bigr),
 \qquad
 \bigl(q^+,\ \text{recursive right subtree of }q\bigr).
\]
Recursion enters the original children at depth $2d+2$. A copy $q^-$
or $q^+$ denotes one active leaf in the indicated region: if this
region begins above depth $2h$, take its all-right descendant leaf.
Other leaves in that region are inactive. An empty original child
also gives an inactive region. The $2n$ active copies have order
\[
 q_1^-<q_1^+<q_2^-<q_2^+<\cdots<q_n^-<q_n^+.
\]
The two copies of one key are adjacent in the active alphabet, though
there can be inactive leaves between them in the full alphabet.

Call the pivots and the two connectors per original node the
\emph{core}. There are $3n$ core nodes. At depth $a$ the number of core
nodes is at most
\begin{equation}\label{eq:core-level}
 2^{\lfloor a/2\rfloor+1}.
\end{equation}
All primary bits of $T$, including the padding, initially point right.

Encode a query to $q$ by the copy on the side opposite to $q$'s current
multi-splay preferred direction. Equivalently, query $q^-$ if the final
switch makes $q$ prefer left, and query $q^+$ if it makes $q$ prefer
right. This defines a leaf-query sequence $Z$ of length $|X|$.

\begin{lemma}\label{lem:core}
Immediately before and after every query, the preferred direction of
$v(q)$ equals the multi-splay direction at $q$. During a query the
changes at even-depth core pivots are precisely the multi-splay ancestor
switches and its final switch. Their depths are $2E_i$.
\end{lemma}
\begin{proof}
At a proper ancestor $q$ of the requested original key, either active
copy lies in the same side as that original key. Thus both processes
make exactly the same ancestor preference changes. At the requested
key itself, the chosen copy is on the opposite side of its pivot and
forces exactly its final toggle. No other original pivot is on the
copy's reference path. The initially equal directions therefore remain
equal by induction. Leaf fictitious bits obey the same induction.
\end{proof}

All other noncompletion emissions in $Z$ are at connector nodes.
Padding paths are all-right paths and are only queried at their
all-right active leaf, so no padding preference ever changes.

\section{The copies are cheap for every competitor}

\begin{lemma}\label{lem:competitor}
For the actual encoded sequence $Z$ on the complete leaf alphabet,
\[
 \OPT(Z)=O(\OPT(X)).
\]
\end{lemma}
\begin{proof}
Replace every node $q$ of an arbitrary competing BST for $X$ by a
two-node gadget: $q^-$ has right child $q^+$, the original left subtree
hangs from $q^-$'s left pointer, and the original right subtree hangs
from $q^+$'s right pointer. A search visits at most two gadget nodes per
original node, and can serve either requested copy.

An original rotation changes only the two adjacent original nodes and
their three hanging subtrees. Its expanded before and after shapes have
at most four gadget nodes and the same ordered hanging subtrees. They
are interconvertible by a constant number of pointer moves and rotations.
The choice of copy may depend on the multi-splay history; the simulation
serves either choice and imposes no restriction on that dependence.

Finally attach the inactive full-alphabet keys as opaque subtrees in
the inorder gaps of the active skeleton. Searches for active copies do
not enter these gap subtrees, and active rotations preserve them. A
comparator may choose this initial full-alphabet tree. Its cost is a
constant multiple of the original execution, without a charge
proportional to the much larger complete reference universe.
\end{proof}

\section{Embedding the secondary ledger and its initialization}

Make a hierarchy $B'$ on $0,\ldots,2h$ by replacing each leaf $d<h$
of $B$ by the two leaves $2d,2d+1$, and replacing leaf $h$ by the
single leaf $2h$. An original internal interval $I=[a,b]$ maps to
\[
 J(I)=
 \begin{cases}
 [2a,2b+1],&b<h,\\
 [2a,2h],&b=h.
 \end{cases}
\]
Its split is at $2c+1$. At the corresponding core owner $v(u)$,
an even emission $2d$ therefore writes exactly the same direction as
the original depth $d$ to $b_{u,I}$.

\begin{lemma}\label{lem:subsequence}
With all-right initialization, $F_{\MS}(X)\le F_{B'}(Z)$.
\end{lemma}
\begin{proof}
Fix one original owner and interval. In the corresponding $B'$ row,
retain only the even core-pivot emissions. By Lemma~\ref{lem:core},
the retained writes, with their chronological order, are precisely the
writes of the original ledger. Removing letters cannot increase the
number of changes in a binary word with a fixed initial letter. Sum
over the injectively corresponding coarse rows of $B'$.
\end{proof}

The nominal complete universe has $2^{2h}$ leaves, so blindly applying
an $O(2^{2h})$ initialization allowance would be useless. The following
support calculation retains the original $O(n)$ allowance.

\begin{lemma}\label{lem:initial}
For the encoded history,
\[
 F_{B'}(Z)\le 2D_{B'}(Z)+2P(Z)+18n.
\]
\end{lemma}
\begin{proof}
The initialized-flip conversion in the Cha-Cha manuscript gives
the more precise intermediate estimate
\begin{equation}\label{eq:precise}
 F_{B'}(Z)\le 2D_{B'}(Z)+3Q_L+2(P(Z)+I_0),
\end{equation}
where $I_0$ counts initial primary changes beyond first-access-free
alternations, and $Q_L$ may be restricted to fragments that ever
receive a left write. Indeed a row receiving only right writes has no
flips and needs no initial-head allowance. Terminal fragments are
already bounded by twice their resets.

Only core primary nodes change, hence $I_0\le3n$. If a fragment
receives a left write, the responsible emission is not the completion
$2h$. It is consequently at a core node. Its owner, at the interval's
left endpoint $a$, is an ancestor of that core node and is itself core.
For an interval split at $c$, $a\le c$, and
\eqref{eq:core-level} bounds the number of such owners by
$2^{\lfloor c/2\rfloor+1}$. Distinct internal intervals of a binary
interval hierarchy have distinct split positions. Therefore
\[
 Q_L\le\sum_{c=0}^{2h-1}2^{\lfloor c/2\rfloor+1}
 =4(2^h-1)=4n.
\]
Substituting these estimates in \eqref{eq:precise} gives the claim.
\end{proof}

\begin{proof}[Proof of Theorem~\ref{thm:endpoint}]
The even core flips are included in all initialized primary flips, so
$S_{\MS}\le P(Z)+3n$. Lemmas~\ref{lem:subsequence} and
\ref{lem:initial} give
\[
 |X|+S_{\MS}+F_{\MS}
 \le |Z|+2D_{B'}(Z)+3P(Z)+21n
 \le3\bigl(|Z|+D_{B'}(Z)+P(Z)\bigr)+21n.
\]
Apply \eqref{eq:lower} and Lemma~\ref{lem:competitor}.
\end{proof}

\section{Scope and the remaining upper-bound question}

\begin{lemma}[The actual switch order]\label{lem:order}
Let $F_{\MS}^{\mathrm{actual}}$ be the same analytical ledger, now fed
the standard multi-splay order. If
$E_i=(d_1<\cdots<d_k)$, that order is
$(d_{k-1},\ldots,d_1,d_k)$. Then
\[
 \bigl|F_{\MS}^{\mathrm{actual}}-F_{\MS}\bigr|
 \le 2\sum_i\max(|E_i|-2,0)\le2S_{\MS}.
\]
Consequently Theorem~\ref{thm:endpoint} also holds with
$F_{\MS}^{\mathrm{actual}}$ in place of $F_{\MS}$.
\end{lemma}
\begin{proof}
Fix one query and restrict its distinct emitted depths to a secondary
interval. A homogeneous left/right block is unchanged. A mixed block
has the form $L^pR^q$ in increasing order. If the interval does not
contain $d_k$, the new block is $R^qL^p$: its internal variation is
unchanged, and each of its two boundary contributions changes by at
most one. If the interval contains $d_k$ and $q=1$, the new block is
still $L^pR$. If it contains $d_k$ and $q\ge2$, the new block is
$R^{q-1}L^pR$: its internal variation increases by one, its last letter
is unchanged, and its entry contribution changes by at most one.
Thus replacing any changed mixed block changes the complete row's
flip count by at most two, including the initialized letter and the
following block if present.

There are exactly $k-1$ secondary intervals meeting both sides of the
emitted set: these are the branching vertices of its induced binary
trie. For $k\ge2$, the parent of the maximum emitted leaf in the
compressed induced trie has only that maximum in its right part.
Its block has $q=1$ and is unchanged. Hence at most $k-2$ mixed blocks
can change. For this query each interval has exactly one owner row;
there is no extra multiplicity from owners. Replace the query blocks
successively in the complete owner-row words and apply the triangle
inequality. This argument does not require a common maximum depth
across queries.
\end{proof}

The fixed-completion version of this order comparison is also recorded
in the companion note \texttt{order-and-potential.tex}. The formulation
above permits the variable final depth of an internal-node query.

\begin{corollary}[Height-optimal references]
Theorem~\ref{thm:endpoint}, including Lemma~\ref{lem:order}, holds for a
node reference on any $n$ keys whose height in vertices is
$h=\lceil\log_2(n+1)\rceil$.
\end{corollary}
\begin{proof}
The copy construction permits absent original subtrees: the
corresponding portions of the complete embedded reference contain only
inactive leaves. There are still at most $3n$ core vertices. The same
level-count estimate gives
$Q_L\le4(2^h-1)<8n+4$, so the initialization allowance remains $O(n)$.
All correspondence, comparator, and order arguments are unchanged.
\end{proof}

\begin{corollary}[Arbitrary initial primary preferences]
The same bounds hold from any valid initial multi-splay preferred-path
partition, with the same all-right initialization of the analytical
secondary ledger.
\end{corollary}
\begin{proof}
Initialize each embedded pivot $v(q)$ to the supplied initial primary
direction at $q$. Keep the connectors and padding initialized right.
The exact pivot correspondence then holds from the start. The
initialized-primary-flip conversion allows an arbitrary initial
direction, with at most one discrepancy per visited core vertex;
thus $I_0\le3n$ still. Padding vertices remain unchanged. All secondary
rows retain their original right initialization, so the subsequence,
row-support, and order proofs apply verbatim.
\end{proof}

For an arbitrary balanced node reference of maximum depth $h-1$, the
same copy embedding and comparator proof still work. The elementary
initialization estimate is then $Q_L=O(n\operatorname{height}(B'))$,
which is $O(n\log h)$ for a balanced depth hierarchy. The present proof
claims the sharper $O(n)$ allowance for a perfect or height-optimal
reference, including the usual recursively median-split reference.
It does not
silently identify $2^h$ with $n$ for a general balanced reference.

The desired next statement would be a physical multi-splay upper bound
of the form
\[
 \operatorname{cost}_{\MS}(X)
 =O\!\left((|X|+S_{\MS}(X)+F_{\MS}(X)+n)
       (1+\log\log(h+2))\right).
\]
This note does not prove that statement. Standard multi-splay uses
bottom-up switches, each implemented by up to three splays and two
root-mark changes. Its splay components split and join. The analytical
top-down ledger above does not itself bound those physical operations.
It only shows that final switches at internal requested nodes need not
be an additional lower-bound obstacle.

\begin{thebibliography}{9}
\bibitem{MST}
C.~C. Wang, J. Derryberry, and D.~D. Sleator.
\emph{$O(\log\log n)$-Competitive Dynamic Binary Search Trees}.
SODA 2006. Sections 2--4, including the final switch convention.
\url{https://www.cs.cmu.edu/~chengwen/paper/MST.pdf}.

\bibitem{TS}
\emph{Cha-Cha: An $O(\log\log\log n)$-Competitive Binary Search
Tree}. Local research manuscript, October 2026.
The dynamic lower-bound theorem, initialized flip conversion, and
their proofs in
\texttt{chacha/research/chacha.tex}.
\end{thebibliography}
\end{document}
